\documentclass[UTF8, 12pt, fontset=fandol, oneside]{ctexart}
\usepackage{researchnotes}

% 保留分册独立编译时的章号前缀。
\renewcommand{\thesection}{10.\arabic{section}}

\title{第十章\quad 函数序列与函数项级数}
\author{}
\date{}

\begin{document}

\maketitle

我们知道数学分析主要是研究函数的一门学科. 到目前为止, 我们经常遇到的函数, 除了一些很特别的函数外(如黎曼函数, 狄利克雷函数等), 大都是一些初等函数. 一个重要的问题是: 是否存在更为广泛的函数类呢? 例如, 是否存在处处连续但处处不可导的函数呢? 再如, 设 $\{n_k\}$ 是一个严格单调上升的正整数序列, 是否存在 $C^\infty(-\infty,+\infty)$ 的函数 $f(x)$, 满足当 $i\in\{n_k\}$ 时, $f^{(i)}(0)=1$, 而当 $i\notin\{n_k\}$ 时, $f^{(i)}(0)=0$ 呢? 仅用前面所学的知识, 似乎很难解决这类问题. 另外, 人们还希望利用简单函数来逼近复杂的函数, 以达到对复杂函数的深入研究. 若要研究上述这些问题, 我们不可避免地要遇到函数序列或函数项级数。

在本章中, 我们将对函数序列及函数项级数进行系统地学习. 本章内容在今后几章和许多后续课程中均起着重要作用。

\section{函数序列与函数项级数的基本问题}

设有一个序列 $\{f_n(x)\}$, 其中序列的第 $n$ 项是一个函数 $f_n(x)$, 而所有 $f_n(x)(n=1,2,\cdots)$ 都有一个共同的定义域 $I_0\subset\mathbb R$. 对于这样的序列, 我们今后称之为定义在 $I_0$ 上的函数序列或称该函数序列在 $I_0$ 上是有定义的. 对于固定的点 $x_0\in I_0$, 该函数序列在点 $x_0$ 的取值组成了一个序列 $\{f_n(x_0)\}$. 若该序列收敛, 我们则称函数序列在 $x_0$ 处收敛, $x_0$ 也称为该函数序列的收敛点. 函数序列 $\{f_n(x)\}$ 的所有收敛点的集合称为它的收敛域. 设 $\{f_n(x)\}$ 的收敛域为 $I\ne\varnothing$, 则对于 $x\in I$, 我们得到 $I$ 上的一个函数 $f(x)$, 它的定义为: $\forall x\in I$, $f(x)=\lim\limits_{n\to\infty}f_n(x)$. $f(x)$ 也称为 $\{f_n(x)\}$ 在 $I$ 上的极限函数。

同样我们可以考虑和式
\[
  \sum_{n=1}^{+\infty}u_n(x),
\]
其中 $\{u_n(x)\}$ 是定义在集合 $I_0$ 上的一个函数序列. 今后我们称该和式为定义在 $I_0$ 上的函数项级数(简称级数) 或称该函数项级数在 $I_0$ 有定义. 记
\[
  S_n(x)=\sum_{k=1}^n u_k(x)\quad(n=1,2,\cdots),
\]
称之为函数项级数 $\sum\limits_{n=1}^{+\infty}u_n(x)$ 的前 $n$ 项部分和, 并称 $\{S_n(x)\}$ 为该函数项级数的部分和序列. 若 $\{S_n(x)\}$ 在 $x_0\in I_0$ 处收敛, 即
\[
  \lim_{n\to\infty}S_n(x_0)=S(x_0),
\]
则称 $\sum\limits_{n=1}^{+\infty}u_n(x)$ 在 $x_0$ 处收敛, $x_0$ 称为 $\sum\limits_{n=1}^{+\infty}u_n(x)$ 的一个收敛点, 其收敛点的全体称为 $\sum\limits_{n=1}^{+\infty}u_n(x)$ 的收敛域. 设 $\sum\limits_{n=1}^{+\infty}u_n(x)$ 的收敛域为 $I$, 则对于 $\forall x\in I$, 记
\[
  \sum_{n=1}^{+\infty}u_n(x)=\lim_{n\to\infty}\sum_{k=1}^n u_k(x)=\lim_{n\to\infty}S_n(x)=S(x),
\]
并称 $S(x)$ 为该函数项级数在 $I$ 上的和函数。

从以上的定义来看, 求函数序列的极限函数 $f(x)$ 与函数项级数的和函数 $S(x)$ 实质上是分别求序列的极限和数项级数的和。

由函数项级数 $\sum\limits_{n=1}^{+\infty}u_n(x)$ 收敛性的定义, 它的敛散性可以转化为函数序列
\[
  \left\{S_n(x)=\sum_{k=1}^n u_k(x)\right\}
\]
的敛散性问题. 反过来, 对于函数序列 $\{f_n(x)\}$, 则它的敛散性也可转化为函数项级数 $\sum\limits_{n=1}^{+\infty}[f_n(x)-f_{n-1}(x)]$ (令 $f_0(x)\equiv0$) 的敛散性问题。

在本章中, 我们所关心的是极限(和)函数与函数序列(函数项级数)中的函数之间具有什么联系? 为此我们先来看以下例子。

\noindent\textbf{例 10.1.1}\quad 函数序列 $f_n(x)=x^n(n=1,2,\cdots)$ 的收敛域是 $(-1,1]$, 其极限函数是
\[
  f(x)=
  \begin{cases}
    0, & |x|<1,\\
    1, & x=1.
  \end{cases}
\]
注意到函数序列 $\{x^n\}$ 中的每一个函数都是连续函数, 其极限函数在 $(-1,1)$ 内连续, 但在 $(-1,1]$ 上却不连续. 同时我们还注意到对于 $\forall n\in\mathbb N$, $f_n(x)$ 均为可导函数, 但其极限函数却在 $x=1$ 处左导数不存在。

\noindent\textbf{例 10.1.2}\quad 设 $f_n(x)=nx(1-x^2)^n(x\in[0,1],n=1,2,\cdots)$, 则
\[
  \lim_{n\to\infty}f_n(x)\equiv0,\quad x\in[0,1],
\]
但
\[
  \lim_{n\to\infty}\int_0^1 f_n(x)\mathrm dx
  =\lim_{n\to\infty}\int_0^1 nx(1-x^2)^n\mathrm dx
  =\frac12
  \ne \int_0^1\lim_{n\to\infty}f_n(x)\mathrm dx=0.
\]

\noindent\textbf{例 10.1.3}\quad 设 $f_n(x)=\dfrac{\sin nx}{\sqrt n}(n=1,2,\cdots)$, 显然
\[
  \lim_{n\to\infty}\frac{\sin nx}{\sqrt n}\equiv0=f(x),\quad x\in(-\infty,+\infty),
\]
但对 $n=1,2,\cdots$, 有
\[
  f_n'(x)=\sqrt n\cos nx.
\]
容易看出 $\{f_n'(x)\}$ 在许多点是不收敛的, 因此该函数序列的导函数序列不收敛于 $f(x)$ 的导数。

\noindent\textbf{例 10.1.4}\quad 设 $f_n(x)=\dfrac{2x}{\pi}\arctan nx(n=1,2,\cdots)$, 显然
\[
  \lim_{n\to\infty}f_n(x)=|x|=f(x),\quad x\in(-\infty,+\infty),
\]
而对 $\forall n\in\mathbb N$, 有
\[
  f_n'(x)=\frac2\pi\arctan nx+\frac{2nx}{\pi(1+n^2x^2)}.
\]
容易看出 $\{f_n'(x)\}$ 处处收敛, 但 $f(x)$ 在 $x=0$ 处不可导。

在讨论函数序列(函数项级数)时, 一个十分重要的问题是研究其极限函数(和函数)能从该函数序列(函数项级数)继承到什么样的性质, 其中主要的有以下三方面的问题(我们主要对函数序列叙述, 读者不难对函数项级数提出相应的问题):

(1) 若 $f_n(x)(n=1,2,\cdots)$ 是区间 $I$ 上的连续函数, 其极限函数 $f(x)$ 是否为 $I$ 上的连续函数?

(2) 若 $f_n(x)(n=1,2,\cdots)$ 是区间 $[a,b]$ 上的可积函数, 其极限函数 $f(x)$ 是否为 $[a,b]$ 上的可积函数? 若 $f(x)\in R[a,b]$, 是否成立
\[
  \lim_{n\to\infty}\int_a^b f_n(x)\mathrm dx=\int_a^b f(x)\mathrm dx?
\]

(3) 若 $f_n(x)(n=1,2,\cdots)$ 是区间 $I$ 上的可导函数, 其极限函数 $f(x)$ 是否为可导函数? 若可导, 是否成立 $\lim\limits_{n\to\infty}f_n'(x)=f'(x)$?

在上面的这些例子中, 我们可以发现: 一般来说, 上述的回答都是否定的. 因此我们必须对所给的函数序列(函数项级数)加上一些条件, 才有可能得到肯定的回答. 在本章以下几节中, 我们将主要研究这些问题。

\section{一致收敛的概念}

在上节中, 我们发现了很多极限函数(和函数)没有保留函数序列(函数项级数)中每个函数的相应性质. 因此一个重要的问题是: 在什么条件下才能保证极限函数能保留相应函数序列中的函数的性质呢? 为此让我们再来考查一下函数序列 $\{f_n(x)=x^n\}$. 在上节中我们已知其收敛域为 $(-1,1]$, 极限函数为
\[
  f(x)=
  \begin{cases}
    0, & |x|<1,\\
    1, & x=1.
  \end{cases}
\]
我们发现 $f(x)$ 在 $(-1,1)$ 内还是保持了函数的连续性, 不仅如此, 当 $x\in(-1,1)$ 时, 还可容易地证明
\[
  \lim_{n\to\infty}(x^n)'=\lim_{n\to\infty}nx^{n-1}=0=\left(\lim_{n\to\infty}x^n\right)'.
\]
为什么极限函数 $f(x)$ 在 $x=1$ 处不是左连续的呢? 我们可以粗略地先来分析一下. 极限函数 $f(x)$ 它在 $x=1$ 左连续与它在 $x=1$ 左边邻域的函数值有关, 而现在 $f(x)$ 的函数值是 $\{x^n\}$ 的极限. 尽管当 $n$ 趋于 $\infty$ 时, $\{x^n\}$ 在 $(1-\delta,1)$ 中都趋于零, 但当 $x$ 靠近 $1$ 时, 使得 $x^n$ 接近零所需要的 $n$ 也越大(见图 10.2.1). 换句话说, 对任给 $\varepsilon>0$, 不管 $n$ 多么大, 对每个固定的 $n$, $|x^n-f(x)|<\varepsilon$ 在 $(1-\delta,1)$ 中不能对所有的 $x$ 同时成立. 因此极限函数就很难保持函数序列中的每一个固定函数所具有的特性了。

另一方面, 在 $(-1,1)$ 中极限函数 $f(x)\equiv0$ 与函数序列 $\{x^n\}$ 收敛到零的过程具有何种联系呢?

我们知道 $f(x)$ 在 $x_0$ 处连续是一个局部的性质. 换句话说, $f(x)$ 在 $x_0$ 处是否连续只与 $f(x)$ 在 $x_0$ 的某个邻域内的函数值有关. 设 $x_0\in(-1,1)$, 可先取定 $\delta_0$ 充分小, 使得 $[x_0-\delta_0,x_0+\delta_0]\subset(-1,1)$. 由 $f(x)$ 在 $x_0$ 处连续的定义, 对于 $\forall\varepsilon>0,\exists\delta>0(\delta<\delta_0)$, 当 $|x-x_0|<\delta$ 时有
\[
  |f(x)-f(x_0)|<\varepsilon.
\]
换句话说, $f(x)$ 在 $x_0$ 处连续本质上是 $f(x)$ 在 $x_0$ 的附近变化不大. 而在 $[x_0-\delta_0,x_0+\delta_0]$ 中, 由于 $(x_0-\delta_0)^n\to0$ 与 $(x_0+\delta_0)^n\to0$, 从而 $\exists N$, 当 $n>N$ 时, 有
\[
  |x^n-f(x)|<\varepsilon
\]
对一切 $x\in[x_0-\delta_0,x_0+\delta_0]$ 成立. 特别地, 对一切 $x\in[x_0-\delta_0,x_0+\delta_0]$, 有
\[
  |x^{N+1}-f(x)|<\varepsilon.
\]
上式说明, 在整个区间 $[x_0-\delta_0,x_0+\delta_0]$ 上, $f(x)$ 与一个连续函数 $x^{N+1}$ 的差的绝对值可以小于任意给定的正数. 由于 $x^{N+1}$ 在 $x_0$ 附近的变化不大, 从而 $f(x)$ 在 $x_0$ 附近的变化也不大, 因此我们从 $n$ 充分大时 $x^n$ 的连续性就可以推出 $f(x)$ 在 $x_0$ 处的连续性。

从上面的分析可以看出 $\{x^n\}$ 在 $(-1,1)$ 中的收敛性与在 $[a,b]\subset(-1,1)$ 中的收敛性具有本质的区别. 对于后者的收敛性, 我们引入以下的重要概念。

\noindent\textbf{定义 10.2.1}\quad 设 $f(x),f_n(x)(n=1,2,\cdots)$ 为定义在 $I\subset\mathbb R$ 上的函数. 若对于 $\forall\varepsilon>0$, 存在 $N\in\mathbb N$, 当 $n>N$ 时, 对一切 $x\in I$, 有
\[
  |f_n(x)-f(x)|<\varepsilon,
\]
则称函数序列 $\{f_n(x)\}$ 在 $I$ 上一致收敛于 $f(x)$, 记为 $f_n(x)\rightrightarrows f(x)(x\in I)$。

显然, 当 $f_n(x)\rightrightarrows f(x)(x\in I)$ 时, $\{f_n(x)\}$ 在 $I$ 上收敛于 $f(x)$。

对于函数项级数, 我们有

\noindent\textbf{定义 10.2.2}\quad 设 $\sum\limits_{n=1}^{+\infty}u_n(x)$ 为定义在 $I\subset\mathbb R$ 上的函数项级数. 若存在 $I$ 上的函数 $S(x)$, 使得 $\sum\limits_{n=1}^{+\infty}u_n(x)$ 的部分和序列 $S_n(x)=\sum\limits_{k=1}^n u_k(x)(n=1,2,\cdots)$ 在 $I$ 上一致收敛到 $S(x)$, 则称 $\sum\limits_{n=1}^{+\infty}u_n(x)$ 在 $I$ 上一致收敛于 $S(x)$。

用 ``$\varepsilon-N$'' 语言来叙述函数项级数的一致收敛性, 我们有: 若对于 $\forall\varepsilon>0$, $\exists N\in\mathbb N$, 当 $n>N$ 时, 对一切 $x\in I$, 有
\[
  |S_n(x)-S(x)|=\left|\sum_{k=n+1}^{+\infty}u_k(x)\right|<\varepsilon,
\]
则称 $\sum\limits_{n=1}^{+\infty}u_n(x)$ 在 $I$ 上一致收敛于 $S(x)$。

下面我们来研究一下一致收敛的几何意义. 设 $\{f_n(x)\}$ 为定义在 $I=[a,b]$ 上的函数序列. 从几何上来看, 对于 $\forall\varepsilon>0$, 则 $y=f(x)\pm\varepsilon(x\in I)$ 的图像与 $x=a,x=b$ 围成了一个带形区域. 当 $\{f_n(x)\}$ 在 $I$ 上一致收敛于 $f(x)$ 时, 则存在 $N\in\mathbb N$, 当 $n>N$ 时, $f_n(x)$ 在 $I$ 上的图像整个地落在该带形区域内(见图 10.2.2)。

\noindent\textbf{例 10.2.1}\quad 试讨论函数序列 $f_n(x)=x^n(n=1,2,\cdots)$ 在

(1) $(-r,r)(0<r<1)$ 内的一致收敛性;

(2) $(-1,1)$ 内的一致收敛性。

\noindent\textbf{解}\quad 我们已经知道 $\lim\limits_{n\to\infty}f_n(x)\equiv0,x\in(-1,1)$。

(1) 对于 $\forall\varepsilon>0$, 由 $r^n\to0(n\to\infty)$, 存在 $N\in\mathbb N$, 当 $n>N$ 时, 有 $|r^n|<\varepsilon$. 因此当 $n>N$ 时, 对于 $\forall x\in(-r,r)$, 都有
\[
  |x^n-0|\leq |r^n|<\varepsilon.
\]
由一致收敛的定义我们有 $x^n\rightrightarrows0(x\in(-r,r))$。

(2) 为了讨论 $f_n(x)$ 在 $(-1,1)$ 内的一致收敛性, 我们首先用肯定的语气来叙述一个函数序列 $\{f_n(x)\}$ 在 $I$ 上不一致收敛于 $f(x)$:

若存在 $\varepsilon_0>0$, 对任意 $N>0$, 存在 $n'>N$ 以及 $x'\in I$, 使得
\[
  |f_{n'}(x')-f(x')|\geq\varepsilon_0,
\]
则 $\{f_n(x)\}$ 在 $I$ 上不一致收敛于 $f(x)$。

现在我们来证明 $\{x^n\}$ 在 $(-1,1)$ 内不一致收敛于零。

事实上, 存在 $\varepsilon_0=\dfrac12$, 对任意 $N>0$, 存在 $n'=N+1>N$, 由于 $\lim\limits_{x\to1}x^{n'}=1$(对固定的 $n'$), 因此存在 $0<x'<1$, 使得
\[
  |(x')^{n'}-0|=(x')^{n'}>1-\frac12=\frac12.
\]
这就证明了 $\{x^n\}$ 在 $(-1,1)$ 内不一致收敛于零。

\noindent\textbf{例 10.2.2}\quad 设函数 $f(x)\in C[0,1]$ 且 $f(1)=0$, 证明:
\[
  x^nf(x)\rightrightarrows0\quad(x\in[0,1]).
\]

\noindent\textbf{证明}\quad 对于 $\forall x\in[0,1]$, 显然有 $\lim\limits_{n\to\infty}x^nf(x)=0$。

现证一致收敛性. 对于 $\forall\varepsilon>0$, 由 $f(x)$ 在 $x=1$ 处左连续, 存在 $\delta>0$, 当 $x\in(1-\delta,1]$ 时, 有
\[
  |f(x)-0|<\varepsilon,
\]
从而对于 $\forall x\in(1-\delta,1]$ 及 $\forall n\in\mathbb N$, 有
\begin{equation}
  |x^nf(x)|\leq |f(x)|<\varepsilon.
  \tag{10.2.1}
\end{equation}
设 $\max\limits_{x\in[0,1]}\{|f(x)|\}=M$, 对于 $\forall x\in[0,1-\delta]$, 则有
\[
  |x^nf(x)|\leq M(1-\delta)^n.
\]
因为 $\lim\limits_{n\to\infty}(1-\delta)^n=0$, 所以 $\exists N\in\mathbb N$, 当 $n>N$ 时, 对于 $\forall x\in[0,1-\delta]$, 有
\[
  |x^nf(x)|<\varepsilon.
\]
注意到当 $x\in(1-\delta,1]$ 时, 对任意的 $n$, (10.2.1) 成立, 从而当 $n>N$ 时, 对于 $\forall x\in[0,1]$, 有 $|x^nf(x)-0|<\varepsilon$. 因此 $x^nf(x)\rightrightarrows0(x\in[0,1])$。

\noindent\textbf{例 10.2.3}\quad 试讨论函数项级数
\[
  \sum_{n=1}^{+\infty}\frac{(2n+1)x}{[1+(n+1)^2x](1+n^2x)}
\]
在

(1) $[a,+\infty)(a>0)$ 上的一致收敛性;

(2) $(0,+\infty)$ 上的一致收敛性。

\noindent\textbf{解}\quad 对于 $\forall n\in\mathbb N$ 和 $\forall x\in(0,+\infty)$, 记
\[
  S_n(x)=\sum_{k=1}^n\frac{(2k+1)x}{[1+(k+1)^2x](1+k^2x)}
  =\frac1{1+x}-\frac1{1+(n+1)^2x},
\]
则有
\[
  \lim_{n\to\infty}S_n(x)=S(x)=\frac1{1+x}.
\]
因此我们有
\[
  \frac1{1+x}=\sum_{n=1}^{+\infty}\frac{(2n+1)x}{[1+(n+1)^2x](1+n^2x)},\quad x\in(0,+\infty).
\]

(1) 当 $x\in[a,+\infty)$ 时, 我们有
\[
  |S_n(x)-S(x)|=\frac1{1+(n+1)^2x}\leq\frac1{1+(n+1)^2a}.
\]
因此对于 $\forall\varepsilon>0$, $\exists N\in\mathbb N$, 当 $n>N$ 时, 有
\[
  0<\frac1{1+(n+1)^2a}<\varepsilon,
\]
从而当 $n>N$ 时, 对于 $\forall x\in[a,+\infty)$, 有
\[
  |S_n(x)-S(x)|<\varepsilon.
\]
这就证明了 $\sum\limits_{n=1}^{+\infty}\dfrac{(2n+1)x}{[1+(n+1)^2x](1+n^2x)}$ 在 $[a,+\infty)(a>0)$ 上一致收敛。

(2) 当 $x\in(0,+\infty)$ 时, 我们观察到: $\exists\varepsilon_0=\dfrac12$, 对于 $\forall N\in\mathbb N$, $\exists n'=N+1>N$, $\exists x'=\dfrac1{(N+1)^2}$, 使得
\[
  |S_{N+1}(x')-S(x')|=\frac1{1+(N+1)^2\dfrac1{(N+1)^2}}=\frac12.
\]
这说明 $\sum\limits_{n=1}^{+\infty}\dfrac{(2n+1)x}{[1+(n+1)^2x](1+n^2x)}$ 在 $(0,+\infty)$ 上不一致收敛。

从定义可以看出, 函数序列的一致收敛具有以下简单性质:

(1) 当函数序列 $\{f_n(x)\}$ 在 $I\subset\mathbb R$ 一致收敛时, 若 $I'\subset I$, 则 $\{f_n(x)\}$ 在 $I'$ 也一致收敛. 另外若 $\{f_n(x)\}$ 分别在 $I_1,I_2\subset\mathbb R$ 上一致收敛, 则它在 $I_1\cup I_2$ 上也一致收敛。

(2) 若函数序列 $\{f_n(x)\}$ 与 $\{g_n(x)\}$ 在 $I\subset\mathbb R$ 上均一致收敛, 则对 $\forall\alpha,\beta\in\mathbb R$, $\{\alpha f_n(x)+\beta g_n(x)\}$ 在 $I$ 上也一致收敛。

容易看出, 对函数项级数而言, 上述两个相应的性质也成立。

人们不禁要问: 如果两个函数序列 $\{f_n(x)\},\{g_n(x)\}$ 都在 $I\subset\mathbb R$ 上一致收敛, 是否 $\{f_n(x)g_n(x)\}$ 在 $I$ 上还一致收敛呢? 一般来说, 答案是否定的. 例如, 取 $I=(0,1)$, $f_n(x)=(1-x)x^n(n=1,2,\cdots)$, 由前面例 10.2.2 知, $\{f_n(x)\}$ 在 $(0,1)$ 内一致收敛. 另外, 令 $g_n(x)=\dfrac1{1-x}(n=1,2,\cdots)$, 则 $\{g_n(x)\}$ 在 $(0,1)$ 内一致收敛. 由于 $f_n(x)g_n(x)=x^n(n=1,2,\cdots)$, 从而 $\{f_n(x)g_n(x)\}$ 在 $(0,1)$ 内不一致收敛。

\noindent\textbf{定义 10.2.3}\quad 设 $\{f_n(x)\}$ 为定义在 $I\subset\mathbb R$ 上的函数序列. 若存在 $M>0$, 使得对于 $\forall n\in\mathbb N$ 和 $\forall x\in I$, 有 $|f_n(x)|\leq M$, 则称 $\{f_n(x)\}$ 在 $I$ 上一致有界。

若函数序列 $\{f_n(x)\}$ 与 $\{g_n(x)\}$ 在 $I\subset\mathbb R$ 都一致收敛, 并且都是一致有界的, 则可以证明 $\{f_n(x)g_n(x)\}$ 在 $I$ 上一致收敛(见本章习题)。

\section{函数序列与函数项级数一致收敛的判别法}

如果一个函数序列 $\{f_n(x)\}$ 一致收敛于 $f(x)$, 此时函数序列中 $f_n(x)$ 是点点收敛于 $f(x)$. 要判别函数序列(函数项级数)是否一致收敛, 一个基本的原则是: 将已有的点点收敛的判别法中的条件换成关于 $x\in I$ 中一致成立, 这样便能得到相应的关于一致收敛的判别法则。

\subsection{柯西准则}

对于函数序列的一致收敛, 我们首先有以下的柯西准则。

\noindent\textbf{定理 10.3.1 (柯西准则)}\quad 设 $\{f_n(x)\}$ 是定义在 $I\subset\mathbb R$ 上的函数序列, 则 $\{f_n(x)\}$ 在 $I$ 上一致收敛的充分必要条件是: 对于 $\forall\varepsilon>0$, 存在 $N\in\mathbb N$, 当 $n,m>N$ 时, 对一切 $x\in I$, 有
\[
  |f_n(x)-f_m(x)|<\varepsilon.
\]

\noindent\textbf{证明}\quad 必要性\quad 设 $f_n(x)\rightrightarrows f(x)(x\in I)$, 则对于 $\forall\varepsilon>0$, 存在 $N\in\mathbb N$, 当 $n>N$ 时, 对于 $\forall x\in I$, 有
\[
  |f_n(x)-f(x)|<\frac\varepsilon2.
\]
因此当 $n,m>N$ 时, 对于 $\forall x\in I$, 有
\[
  |f_n(x)-f_m(x)|\leq |f_n(x)-f(x)|+|f_m(x)-f(x)|<\varepsilon.
\]
必要性得证。

充分性\quad 设 $\{f_n(x)\}$ 满足: 对于 $\forall\varepsilon>0$, 存在 $N\in\mathbb N$, 当 $n,m>N$ 时, 对一切 $x\in I$, 有
\[
  |f_n(x)-f_m(x)|<\frac\varepsilon2.
\]
特别地, 上式对每个固定的点 $x\in I$ 也成立, 从而序列 $\{f_n(x)\}$ 收敛, 设其极限为 $f(x)$, 因此我们在区间 $I$ 上可得到 $\{f_n(x)\}$ 的极限函数 $f(x)$. 在
\[
  |f_n(x)-f_m(x)|<\frac\varepsilon2
\]
中令 $m\to\infty$, 我们有
\[
  |f_n(x)-f(x)|\leq\frac\varepsilon2<\varepsilon
\]
对一切 $n>N$ 及 $x\in I$ 成立. 这就证明了 $f_n(x)\rightrightarrows f(x)(x\in I)$. 证毕。

对于函数项级数, 我们有

\noindent\textbf{定理 10.3.1$'$ (柯西准则)}\quad 设 $\sum\limits_{n=1}^{+\infty}u_n(x)$ 是定义在 $I\subset\mathbb R$ 上的函数项级数, 则 $\sum\limits_{n=1}^{+\infty}u_n(x)$ 在 $I$ 上一致收敛的充分必要条件是: 对于 $\forall\varepsilon>0$, 存在 $N\in\mathbb N$, 当 $n>m>N$ 时, 对一切 $x\in I$, 有
\[
  \left|\sum_{k=m+1}^n u_k(x)\right|<\varepsilon.
\]
定理 10.3.1$'$ 的证明留给读者。

从定理 10.3.1$'$ 可以推出下面有用的推论。

\noindent\textbf{推论}\quad 设 $\sum\limits_{n=1}^{+\infty}u_n(x)$ 是定义在 $I\subset\mathbb R$ 上的函数项级数, 并且它在 $I$ 上一致收敛, 则 $u_n(x)\rightrightarrows0(x\in I)$。

\noindent\textbf{例 10.3.1}\quad 设函数 $f_n(x)(n=1,2,\cdots)$ 在闭区间 $[a,b]$ 上连续, 并且 $\{f_n(x)\}$ 在开区间 $(a,b)$ 内一致收敛. 证明 $\{f_n(x)\}$ 在 $[a,b]$ 上一致收敛。

\noindent\textbf{证明}\quad 由于 $\{f_n(x)\}$ 在 $(a,b)$ 内一致收敛, 由柯西准则, 对于 $\forall\varepsilon>0$, 存在 $N\in\mathbb N$, 当 $n,m>N$ 时, 对于 $\forall x\in(a,b)$, 有
\[
  |f_n(x)-f_m(x)|<\frac\varepsilon2.
\]
对每个 $n,m>N$, 由于 $f_n(x),f_m(x)$ 在 $[a,b]$ 上连续, 因此在上述不等式中令 $x\to a+0$ 和 $x\to b-0$ 得
\[
  |f_n(a)-f_m(a)|\leq\frac\varepsilon2<\varepsilon
\]
和
\[
  |f_n(b)-f_m(b)|\leq\frac\varepsilon2<\varepsilon.
\]
这说明当 $n,m>N$ 时, 对一切 $x\in[a,b]$, 有
\[
  |f_n(x)-f_m(x)|<\varepsilon.
\]
再由柯西准则知 $\{f_n(x)\}$ 在 $[a,b]$ 上一致收敛。

读者同样可以证明以下结果:

\noindent\textbf{例 10.3.1$'$}\quad 设函数 $u_n(x)(n=1,2,\cdots)$ 在闭区间 $[a,b]$ 上连续, 且 $\sum\limits_{n=1}^{+\infty}u_n(x)$ 在开区间 $(a,b)$ 内一致收敛, 则 $\sum\limits_{n=1}^{+\infty}u_n(x)$ 在 $[a,b]$ 上一致收敛。

\noindent\textbf{例 10.3.2}\quad 试讨论函数项级数 $\sum\limits_{n=1}^{+\infty}\dfrac{(-1)^n}{1+nx}$ 在

(1) $[a,+\infty)(a>0)$ 上的一致收敛性;

(2) $(0,+\infty)$ 上的一致收敛性。

\noindent\textbf{证明}\quad 显然, 对任意固定的 $x\in(0,+\infty)$, 该级数为一个莱布尼茨交错级数, 因此是收敛的. 对于 $\forall n\in\mathbb N$, 记
\[
  S_n(x)=\sum_{k=1}^n\frac{(-1)^k}{1+kx}\quad\text{及}\quad S(x)=\sum_{n=1}^{+\infty}\frac{(-1)^n}{1+nx}.
\]

(1) 当 $x\in[a,+\infty)$ 时, 由交错级数余项的性质得(见定理 9.3.1 推论后面的注)
\[
  |S_n(x)-S(x)|=\left|\sum_{k=n+1}^{+\infty}\frac{(-1)^k}{1+kx}\right|\leq\frac1{1+(n+1)x}\leq\frac1{1+(n+1)a}.
\]
由于对于 $\forall\varepsilon>0$, $\exists N\in\mathbb N$, 当 $n>N$ 时, 有
\[
  \frac1{1+(n+1)a}<\varepsilon,
\]
因此当 $n>N$ 时, 对于 $\forall x\in[a,+\infty)$, 有
\[
  |S_n(x)-S(x)|<\varepsilon.
\]
这就证明了 $\sum\limits_{n=1}^{+\infty}\dfrac{(-1)^n}{1+nx}$ 在 $[a,+\infty)$ 上一致收敛。

(2) 注意到 $u_n(x)=\dfrac{(-1)^n}{1+nx}$ 在 $[0,+\infty)$ 上连续, 若 $\sum\limits_{n=1}^{+\infty}u_n(x)$ 在 $(0,+\infty)$ 上一致收敛, 由例 10.3.1$'$ 知它在 $x=0$ 必收敛. 由于 $\sum\limits_{n=1}^{+\infty}u_n(0)=\sum\limits_{n=1}^{+\infty}(-1)^n$ 是发散的, 因此 $\sum\limits_{n=1}^{+\infty}u_n(x)$ 在 $(0,+\infty)$ 上不一致收敛。

\noindent\textbf{注}\quad (2) 也可以用以下的方法证明. 由定理 10.3.1$'$ 的推论可知, 我们只要证明 $u_n(x)=\dfrac{(-1)^n}{1+nx}$ 在 $(0,+\infty)$ 上不一致收敛到零即可. 事实上, $\exists\varepsilon_0=\dfrac12$, 对于 $\forall N\in\mathbb N$, $\exists N+1>N$ 及 $x'=\dfrac1{N+1}\in(0,+\infty)$, 使得
\[
  \left|\frac{(-1)^{N+1}}{1+(N+1)x'}\right|=\frac12=\varepsilon_0.
\]
这就证明了 $\sum\limits_{n=1}^{+\infty}\dfrac{(-1)^n}{1+nx}$ 在 $(0,+\infty)$ 上不一致收敛。

\subsection{一致收敛的判别法}

柯西准则是一个在理论上非常重要的定理, 但有时在具体应用上不是很方便. 以下的判别法在实用上是比较方便的。

\noindent\textbf{定理 10.3.2}\quad 设函数序列 $\{f_n(x)\}$ 在集合 $I\subset\mathbb R$ 上有定义, 则 $f_n(x)\rightrightarrows f(x)(x\in I)$ 的充分必要条件是 $\lim\limits_{n\to\infty}\sup\limits_{x\in I}\{|f_n(x)-f(x)|\}=0$。

\noindent\textbf{证明}\quad 必要性\quad 设 $f_n(x)\rightrightarrows f(x)(x\in I)$, 则对于 $\forall\varepsilon>0$, 存在 $N\in\mathbb N$, 当 $n>N$ 时, 对一切 $x\in I$, 有 $|f_n(x)-f(x)|<\dfrac\varepsilon2$. 因此我们有
\[
  0\leq \sup_{x\in I}\{|f_n(x)-f(x)|\}\leq\frac\varepsilon2<\varepsilon,
\]
此即 $\lim\limits_{n\to\infty}\sup\limits_{x\in I}\{|f_n(x)-f(x)|\}=0$. 必要性得证。

充分性\quad 由于对于 $\forall\varepsilon>0$, 存在 $N\in\mathbb N$, 当 $n>N$ 时, 有
\[
  \sup_{x\in I}\{|f_n(x)-f(x)|\}<\varepsilon,
\]
因此, 当 $n>N$ 时, 对一切 $x\in I$, 有
\[
  |f_n(x)-f(x)|\leq \sup_{x\in I}\{|f_n(x)-f(x)|\}<\varepsilon,
\]
此即 $f_n(x)\rightrightarrows f(x)(x\in I)$. 证毕。

以上定理告诉我们, 当我们知道一个函数序列的极限函数时, 判断它的一致收敛性的问题可转化为求上确界以及讨论序列的收敛性问题, 以上的判别法被人们称为最值判别法。

\noindent\textbf{例 10.3.3}\quad 试讨论 $f_n(x)=\dfrac{x}{1+n^2x^2}(x\in(-\infty,+\infty), n=1,2,\cdots)$ 的一致收敛性。

\noindent\textbf{解}\quad 对于 $\forall x\in(-\infty,+\infty)$, 不难看出 $f_n(x)\to0(n\to\infty)$. 对每个固定的 $n$, 我们来求 $|f_n(x)|$ 在 $(-\infty,+\infty)$ 上的上确界. 由于 $f_n(x)$ 是奇函数且 $f_n(x)\geq0(x\geq0)$, 因此只要计算 $f_n(x)$ 在 $x\geq0$ 时的上确界即可。

由
\[
  f_n'(x)=\frac{(1+n^2x^2)-x\cdot2n^2x}{(1+n^2x^2)^2}
  =\frac{1-n^2x^2}{(1+n^2x^2)^2}=0
\]
得 $x=\pm\dfrac1n$, 注意到 $f(0)=0$ 和 $\lim\limits_{x\to+\infty}f_n(x)=0$, 知 $f_n(x)$ 在 $\dfrac1n$ 取到最大值 $\dfrac1{2n}$。

因为
\[
  \lim_{n\to\infty}\sup_{x\in(-\infty,+\infty)}|f_n(x)-0|
  =\lim_{n\to\infty}\frac1{2n}=0,
\]
所以
\[
  \frac{x}{1+n^2x^2}\rightrightarrows0\quad(x\in(-\infty,+\infty)).
\]

\noindent\textbf{例 10.3.4}\quad 证明函数序列
\[
  f_n(x)=n^2\left(e^{\frac1{nx}}-1\right)\sin\frac1{nx}\quad(n=1,2,\cdots)
\]
在 $[a,+\infty)(a>0)$ 上一致收敛。

\noindent\textbf{证明}\quad 对每个 $x\in[a,+\infty)$, 我们有
\[
  \lim_{n\to\infty}f_n(x)
  =\lim_{n\to\infty}\frac{\sin\frac1{nx}}{\frac1{nx}}\cdot
  \frac{e^{\frac1{nx}}-1}{\frac1{nx}}\cdot\frac1{x^2}
  =\frac1{x^2}.
\]
注意到当 $t\to0$ 时, 有
\[
  e^t-1\sim t+o(t),\qquad \sin t\sim t+o(t),
\]
从而当 $x\in[a,+\infty)$ 且 $n\to\infty$ 时, 有
\[
\begin{aligned}
  \left|n^2\left(e^{\frac1{nx}}-1\right)\sin\frac1{nx}-\frac1{x^2}\right|
  &=n^2\left|\left(e^{\frac1{nx}}-1\right)\sin\frac1{nx}-\frac1{n^2x^2}\right|\\
  &=n^2\left|\left[\frac1{nx}+o\left(\frac1{nx}\right)\right]
  \left[\frac1{nx}+o\left(\frac1{nx}\right)\right]-\frac1{n^2x^2}\right|\\
  &=n^2o\left(\frac1{n^2x^2}\right).
\end{aligned}
\]
因此我们有
\[
  \sup_{x\in[a,+\infty)}
  \left|n^2\left(e^{\frac1{nx}}-1\right)\sin\frac1{nx}-\frac1{x^2}\right|
  \leq n^2o\left(\frac1{n^2a^2}\right)\to0\quad(n\to\infty).
\]
这就证明了 $f_n(x)\rightrightarrows\dfrac1{x^2}(x\in[a,+\infty))$。

对于函数项级数来说, 我们也可以对其部分和序列利用定理 10.3.2 来判别其收敛性. 如我们可以用定理 10.3.2 来讨论 $\left\{\dfrac{1-x^{n+1}}{1-x}\right\}$, 从而可以解决 $\sum\limits_{n=1}^{+\infty}x^n$ 在 $(-1,1)$ 内的一致收敛性问题. 但在一般情形, 要求出一个函数项级数部分和序列的通项表达式不是一件容易的事情, 因此就很难用定理 10.3.2 来判别其一致收敛性。

对于函数项级数, 我们有以下的概念。

\noindent\textbf{定义 10.3.1}\quad 设 $\sum\limits_{n=1}^{+\infty}u_n(x)$ 是定义在 $I$ 上的函数项级数, 并且 $\sum\limits_{n=1}^{+\infty}|u_n(x)|$ 在 $I$ 上一致收敛, 则称 $\sum\limits_{n=1}^{+\infty}u_n(x)$ 在 $I$ 上绝对一致收敛。

利用柯西准则, 读者易于证明, 若 $\sum\limits_{n=1}^{+\infty}u_n(x)$ 在 $I$ 上绝对一致收敛, 则它在 $I$ 上一致收敛。

对于函数项级数, 以下的判别法是经常用到的。

\noindent\textbf{定理 10.3.3 (魏尔斯特拉斯 M 判别法)}\quad 设函数项级数 $\sum\limits_{n=1}^{+\infty}u_n(x)$ 在 $I\subset\mathbb R$ 有定义. 若存在正数序列 $\{M_n\}$, 使得对每个 $n\in\mathbb N$, 对一切 $x\in I$, 有
\[
  |u_n(x)|\leq M_n,
\]
并且 $\sum\limits_{n=1}^{+\infty}M_n$ 收敛, 则 $\sum\limits_{n=1}^{+\infty}u_n(x)$ 在 $I$ 上绝对一致收敛。

\noindent\textbf{证明}\quad 由于 $\sum\limits_{n=1}^{+\infty}M_n$ 收敛, 根据数项级数收敛的柯西准则, 对于 $\forall\varepsilon>0$, 存在 $N\in\mathbb N$, 当 $n>m>N$ 时, 有
\[
  \sum_{k=m+1}^n M_k<\varepsilon.
\]
从而对一切 $x\in I$, 有
\[
  \sum_{k=m+1}^n |u_k(x)|\leq\sum_{k=m+1}^n M_k<\varepsilon,
\]
再由函数项级数一致收敛的柯西准则知 $\sum\limits_{n=1}^{+\infty}u_n(x)$ 在 $I$ 上绝对一致收敛. 证毕。

\noindent\textbf{例 10.3.5}\quad 证明 $\sum\limits_{n=1}^{+\infty}x^ke^{-nx}(k>1)$ 在 $[0,+\infty)$ 上一致收敛。

\noindent\textbf{证明}\quad 对每个正整数 $n$, 当 $x\in(0,+\infty)$ 时, $u_n(x)=x^ke^{-nx}>0$. 注意到 $u_n(0)=0$ 和 $\lim\limits_{x\to+\infty}x^ke^{-nx}=0(n=1,2,\cdots)$, 因此 $u_n(x)$ 在 $(0,+\infty)$ 上取到最大值. 由
\[
  u_n'(x)=kx^{k-1}e^{-nx}-nx^ke^{-nx}=0
\]
得 $x=0$ 和 $x=\dfrac{k}{n}$. 因此 $u_n(x)$ 在 $x=\dfrac{k}{n}$ 取最大值 $e^{-k}\dfrac{k^k}{n^k}$。

由于 $\sum\limits_{n=1}^{+\infty}e^{-k}\dfrac{k^k}{n^k}$ 当 $k>1$ 时收敛, 由魏尔斯特拉斯 M 判别法知 $\sum\limits_{n=1}^{+\infty}x^ke^{-nx}$ 在 $[0,+\infty)$ 上一致收敛。

对于不是绝对一致收敛的函数项级数, 我们有下面的狄利克雷判别法与阿贝尔判别法。

\noindent\textbf{定理 10.3.4 (狄利克雷判别法)}\quad 设函数序列 $u_n(x),v_n(x)(n=1,2,\cdots)$ 在 $I\subset\mathbb R$ 上有定义, 并且满足以下条件:

(1) $\sum\limits_{n=1}^{+\infty}u_n(x)$ 的部分和序列在 $I$ 上一致有界, 即存在 $M>0$, 使得对于 $\forall n\geq1$ 及 $\forall x\in I$, 有
\[
  |S_n(x)|=\left|\sum_{k=1}^n u_k(x)\right|\leq M;
\]

(2) 对每个 $x\in I$, $\{v_n(x)\}$ 关于 $n$ 是单调的, 且 $v_n(x)\rightrightarrows0(x\in I)$,

则 $\sum\limits_{n=1}^{+\infty}u_n(x)v_n(x)$ 在 $I$ 上一致收敛。

\noindent\textbf{证明}\quad 由条件 (1), 当 $n>m\geq1$ 时, 对一切 $x\in I$, 我们有
\[
  \left|\sum_{k=m+1}^n u_k(x)\right|=|S_n(x)-S_m(x)|\leq2M.
\]
由 $v_n(x)\rightrightarrows0(x\in I)$, 对于 $\forall\varepsilon>0$, $\exists N\in\mathbb N$, 当 $n>N$ 时, 对于 $\forall x\in I$ 有
\[
  |v_n(x)|<\frac{\varepsilon}{4M}.
\]
注意到对每个 $x\in I$, $\{v_k(x)\}$ 关于 $n$ 是单调的, 如同数项级数的狄利克雷判别法的证明, 利用阿贝尔变换 (引理 7.5.4) 知, 当 $n>m>N$ 时, 对 $\forall x\in I$, 有
\[
  \left|\sum_{k=m+1}^n u_k(x)v_k(x)\right|
  \leq M(|v_{m+1}(x)|+2|v_n(x)|)<\varepsilon.
\]
由柯西准则知 $\sum\limits_{n=1}^{+\infty}u_n(x)v_n(x)$ 在 $I$ 上一致收敛. 证毕。

\noindent\textbf{定理 10.3.5 (阿贝尔判别法)}\quad 设函数序列 $u_n(x),v_n(x)(n=1,2,\cdots)$ 在 $I\subset\mathbb R$ 上有定义, 并且满足以下条件:

(1) $\sum\limits_{n=1}^{+\infty}u_n(x)$ 在 $I$ 上一致收敛;

(2) 对固定的 $x\in I$, $\{v_n(x)\}$ 关于 $n$ 是单调函数, 且 $\{v_n(x)\}$ 在 $I$ 上一致有界,

则 $\sum\limits_{n=1}^{+\infty}u_n(x)v_n(x)$ 在 $I$ 上一致收敛。

\noindent\textbf{证明}\quad 由条件 (2), 存在 $M>0$, 使得对于 $\forall n\geq1$ 及 $\forall x\in I$, 有
\[
  |v_n(x)|\leq M.
\]
再由 (1) 及柯西准则, 对于 $\forall\varepsilon>0$, $\exists N\in\mathbb N$, 当 $n>m>N$ 时, 对一切 $x\in I$, 有
\[
  \left|\sum_{k=m+1}^n u_k(x)\right|\leq\frac{\varepsilon}{3M}.
\]
由于 $\{v_n(x)\}$ 关于 $n$ 是单调的, 由阿贝尔变换 (引理 7.5.4), 对任意的 $n>m>N$ 及 $\forall x\in I$, 有
\[
  \left|\sum_{k=m+1}^n u_k(x)v_n(x)\right|
  \leq\frac{\varepsilon}{3M}(|v_{m+1}(x)|+2|v_n(x)|)<\varepsilon.
\]
再由柯西准则知 $\sum\limits_{n=1}^{+\infty}u_n(x)v_n(x)$ 在 $I$ 上一致收敛. 证毕。

回忆一下, 在数项级数 $\sum\limits_{n=1}^{+\infty}a_nb_n$ 的阿贝尔判别法的证明中, 我们直接设单调序列 $\{b_n\}$ 的极限为 $b$, 然后对 $\sum\limits_{n=1}^{+\infty}a_n(b_n-b)$ 应用狄利克雷判别法即可得到证明. 对函数项级数 $\sum\limits_{n=1}^{+\infty}u_n(x)v_n(x)$ 而言, 对每点 $x\in I$, $\lim\limits_{n\to\infty}v_n(x)$ 存在, 因此我们也有极限函数 $v(x)$. 为了要应用函数项级数的狄利克雷判别法, 我们必须要证明 $v_n(x)\rightrightarrows v(x)(x\in I)$. 这个结论是否正确呢? 一般来说此结论是不成立, 例如 $\{x^n\}$ 在 $[0,1]$ 上关于 $n$ 单调并且一致有界, 但 $\{x^n\}$ 在该区间上并不一致收敛。

\noindent\textbf{例 10.3.6}\quad 证明函数项级数 $\sum\limits_{n=1}^{+\infty}\dfrac{(-1)^n}{n+x}$ 在

(1) $[0,+\infty)$ 上一致收敛;

(2) 任何不含负整数点的闭区间 $I$ 上一致收敛;

(3) 任何区间 $I$ 上不绝对收敛。

\noindent\textbf{证明}\quad (1) 对任意的正整数 $n$, 令 $u_n(x)=(-1)^n$, 则对于 $x\in[0,+\infty)$, 有
\[
  \left|\sum_{k=1}^n u_k(x)\right|\leq1.
\]
令 $v_n(x)=\dfrac1{n+x}(n=1,2,\cdots)$, 显然 $v_n(x)$ 对一切 $x\in[0,+\infty)$ 均是关于 $n$ 单调下降的函数, 且当 $x\in[0,+\infty)$ 时, 有
\[
  |v_n(x)|\leq\frac1n\to0\quad(n\to\infty).
\]
因此对于 $x\in[0,+\infty)$, 有
\[
  v_n(x)\rightrightarrows0\quad(x\in[0,+\infty)).
\]
根据狄利克雷判别法, $\sum\limits_{n=1}^{+\infty}\dfrac{(-1)^n}{n+x}$ 在 $[0,+\infty)$ 上一致收敛。

(2) 若闭区间 $I\subset[0,+\infty)$, 显然结论成立. 不妨设 $I\subset(-K,-K+1), K\in\mathbb N$. 如同 (1), 我们容易验证: 对于 $\forall n\in\mathbb N$ 及 $\forall x\in I$, $\left|\sum\limits_{k=1}^n u_k(x)\right|\leq1$. 当 $n>K$ 时, $v_n(x)$ 对于 $x\in I$ 均是关于 $n$ 单调上升的且 $v_n(x)\rightrightarrows0(x\in I)$. 因此由狄利克雷判别法, $\sum\limits_{n=1}^{+\infty}\dfrac{(-1)^n}{n+x}$ 在 $I$ 上一致收敛。

(3) $\sum\limits_{n=1}^{+\infty}\dfrac1{n+x}$ 对所有非负整数的点 $x$, 当 $n$ 充分大时, 均是正项级数, 且 $\dfrac1{n+x}\sim\dfrac1n(n\to\infty)$. 由于 $\sum\limits_{n=1}^{+\infty}\dfrac1n$ 发散, 因此 $\sum\limits_{n=1}^{+\infty}\dfrac1{n+x}$ 也发散. 因此 $\sum\limits_{n=1}^{+\infty}\dfrac{(-1)^n}{n+x}$ 在任何区间都不是绝对收敛。

\noindent\textbf{例 10.3.7}\quad 证明函数项级数 $\sum\limits_{n=1}^{+\infty}\dfrac{a_n}{n^x}$ 在 $[1,+\infty)$ 上一致收敛的充分必要条件是 $\sum\limits_{n=1}^{+\infty}\dfrac{a_n}{n}$ 收敛。

\noindent\textbf{证明}\quad 必要性是显然的, 现证充分性。

设 $\sum\limits_{n=1}^{+\infty}\dfrac{a_n}{n}$ 收敛, 在
\[
  \sum_{n=1}^{+\infty}\frac{a_n}{n^x}
  =\sum_{n=1}^{+\infty}\frac{a_n}{nn^{x-1}}
\]
中, 对于 $n=1,2,\cdots$, 我们令
\[
  u_n(x)=\frac{a_n}{n},\qquad v_n(x)=\frac1{n^{x-1}}.
\]
因为 $\sum\limits_{n=1}^{+\infty}u_n(x)=\sum\limits_{n=1}^{+\infty}\dfrac{a_n}{n}$ 收敛, 故它在 $[1,+\infty)$ 上一致收敛. 由于 $\{v_n(x)\}$ 对固定的 $x$ 是 $n$ 的单调函数, 且对于 $\forall x\in[1,+\infty)$ 及 $\forall n\in\mathbb N$, 有
\[
  |v_n(x)|=\left|\frac1{n^{x-1}}\right|\leq1,
\]
根据阿贝尔判别法, $\sum\limits_{n=1}^{+\infty}\dfrac{a_n}{n^x}$ 在 $[1,+\infty)$ 上一致收敛。

\noindent\textbf{例 10.3.8}\quad 试讨论函数项级数 $\sum\limits_{n=1}^{+\infty}(-1)^n(1-x)x^n$ 在区间 $[0,1]$ 上的

(1) 一致收敛性;

(2) 绝对收敛性;

(3) 绝对一致收敛性。

\noindent\textbf{解}\quad (1) 取 $u_n(x)=(-1)^n$, $v_n(x)=(1-x)x^n(n=1,2,\cdots)$, 则 $\left\{\sum\limits_{k=1}^n u_k(x)\right\}$ 在 $[0,1]$ 上一致有界, 而 $v_n(x)$ 对固定的 $x$ 是 $n$ 的单调下降函数, 而且 $v_n(x)\rightrightarrows0(x\in[0,1])$ (见例 10.2.2). 根据狄利克雷判别法, $\sum\limits_{n=1}^{+\infty}(-1)^n(1-x)x^n$ 在 $[0,1]$ 上一致收敛。

(2) 对于 $n=1,2,\cdots$, 记
\[
  S_n(x)=\sum_{k=1}^n(1-x)x^n=
  \begin{cases}
    x-x^{n+1}, & x\in[0,1),\\
    0, & x=1.
  \end{cases}
\]
注意当 $x\in[0,1)$ 时,
\[
  \lim_{n\to\infty}S_n(x)=S(x)=x.
\]
因此当 $x\in[0,1]$ 时, $\sum\limits_{n=1}^{+\infty}|(-1)^n(1-x)x^n|$ 收敛, 从而原级数在 $[0,1]$ 上绝对收敛。

(3) 继续使用 (2) 中的记号. 由于在 $[0,1)$ 中 $|S_n(x)-S(x)|=x^{n+1}$, 而我们已经知道 $\{x^n\}$ 在 $[0,1)$ 上不一致收敛, 因此 $\sum\limits_{n=1}^{+\infty}(-1)^n(1-x)x^n$ 在 $[0,1]$ 不是绝对一致收敛的。

此例说明了即使一个函数项级数在一个区间上一致收敛并且绝对收敛, 但该级数在该区间上不一定绝对一致收敛。

\section{一致收敛的函数序列和函数项级数}

在本节中, 我们来讨论一致收敛的函数序列(函数项级数)的极限函数(和函数)的性质, 并且回答 §10.1 中提出的问题。

\subsection{极限函数的连续性}

\noindent\textbf{定理 10.4.1}\quad 设函数 $f_n(x)\in C[a,b](n=1,2,\cdots)$, 且 $f_n(x)\rightrightarrows f(x)(x\in[a,b])$, 则 $f(x)\in C[a,b]$。

\noindent\textbf{证明}\quad 任取 $x_0\in[a,b]$, 我们证 $f(x)$ 在 $x_0$ 处连续. 对于 $\forall\varepsilon>0$, 由于 $f_n(x)\rightrightarrows f(x)(x\in[a,b])$, 所以 $\exists N\in\mathbb N$, 当 $n>N$ 时, 对一切 $x\in[a,b]$, 有
\[
  |f_n(x)-f(x)|<\frac\varepsilon3.
\]
取定 $n_0>N$(例如取 $n_0=N+1$), 由于 $f_{n_0}(x)$ 在 $x_0\in[a,b]$ 处连续, 从而存在 $\delta>0$, 使得当 $x\in U(x_0,\delta)\cap[a,b]$ 时, 有
\[
  |f_{n_0}(x)-f_{n_0}(x_0)|<\frac\varepsilon3.
\]
因此当 $x\in U(x_0,\delta)\cap[a,b]$ 时, 有
\[
\begin{aligned}
  |f(x)-f(x_0)|
  &\leq |f(x)-f_{n_0}(x)|+|f_{n_0}(x)-f_{n_0}(x_0)|+|f_{n_0}(x_0)-f(x_0)|\\
  &<\varepsilon.
\end{aligned}
\]
证毕。

\noindent\textbf{注}\quad 仿照定理 10.4.1 的证明方法我们可以证明以下结果: 设 $\{f_n(x)\}$ 是定义在 $[a,b]\setminus\{x_0\}$ 上的函数序列, 其中 $x_0\in[a,b]$. 如果 $f_n(x)\rightrightarrows f(x), (x\in[a,b]\setminus\{x_0\})$, 并且对每个 $n\geq1$, 有 $\lim\limits_{x\to x_0}f_n(x)=c_n$, 则 $\lim\limits_{n\to+\infty}c_n$ 存在, 并且有
\[
  \lim_{x\to x_0}f(x)=\lim_{n\to+\infty}c_n.
\]

对函数项级数, 用完全类似于定理 10.4.1 的证明方法, 可以证明

\noindent\textbf{定理 10.4.1$'$}\quad 设函数 $u_n(x)\in C[a,b](n=1,2,\cdots)$, 且 $\sum\limits_{n=1}^{+\infty}u_n(x)$ 在 $[a,b]$ 上一致收敛, 则 $\sum\limits_{n=1}^{+\infty}u_n(x)$ 的和函数在 $[a,b]$ 上连续。

\noindent\textbf{注}\quad 用类似的方法我们还可以证明以下结论: 若 $\sum\limits_{n=1}^{+\infty}u_n(x)$ 在 $I\setminus\{x_0\}$ 上一致收敛, 且对每个 $n\geq1$, 有 $\lim\limits_{x\to x_0}u_n(x)=a_n$, 则 $\sum\limits_{n=1}^{+\infty}a_n$ 收敛, 且
\[
  \lim_{x\to x_0}\sum_{n=1}^{+\infty}u_n(x)
  =\sum_{n=1}^{+\infty}\left[\lim_{x\to x_0}u_n(x)\right]
  =\sum_{n=1}^{+\infty}a_n.
\]

从定理 10.4.1 我们可知, 当连续的函数序列 $\{f_n(x)\}$(即对每个 $n$, $f_n(x)$ 连续) 一致收敛时, 有以下极限顺序的交换性, 即
\[
  \lim_{x\to x_0}\lim_{n\to\infty}f_n(x)
  =\lim_{n\to\infty}\lim_{x\to x_0}f_n(x).
\]
从定理 10.4.1$'$ 可知, 当连续的函数项级数 $\sum\limits_{n=1}^{+\infty}u_n(x)$(即对每个 $n$, $u_n(x)$ 连续) 一致收敛时, 有以下极限号与求和号的可交换性
\[
  \lim_{x\to x_0}\sum_{n=1}^{+\infty}u_n(x)
  =\sum_{n=1}^{+\infty}\lim_{x\to x_0}u_n(x).
\]
我们曾经多次指出, 函数连续是一个局部性的概念. 换句话说, 函数 $f(x)$ 在 $x=x_0$ 处连续仅与 $f(x)$ 在 $x_0$ 附近的函数值有关. 因此若一个连续函数序列在 $x_0$ 的一个邻域内一致收敛到 $f(x)$(此时也称它在 $x_0$ 处是局部一致收敛的), 则 $f(x)$ 在 $x_0$ 处连续. 在数学的一些分支中, 以下内闭一致收敛的概念经常用到。

\noindent\textbf{定义 10.4.1}\quad 设函数序列 $\{f_n(x)\}$ (函数项级数 $\sum\limits_{n=1}^{+\infty}u_n(x)$) 在区间 $I$ 上有定义, 若对任何的闭区间 $[a,b]\subset I$, $\{f_n(x)\}$ ($\sum\limits_{n=1}^{+\infty}u_n(x)$) 在 $[a,b]$ 上一致收敛, 则称 $\{f_n(x)\}$ ($\sum\limits_{n=1}^{+\infty}u_n(x)$) 在 $I$ 内闭一致收敛。

读者不难证明, 对于开区间 $I$ 来说, 内闭一致收敛等价于局部一致收敛. 因此我们有

\noindent\textbf{推论}\quad (1) 设函数 $f_n(x)\in C(a,b)(n=1,2,\cdots)$ 且 $\{f_n(x)\}$ 在 $(a,b)$ 内闭一致收敛于 $f(x)$, 则 $f(x)\in C(a,b)$;

(2) 设函数 $u_n(x)\in C(a,b)$, 且 $\sum\limits_{n=1}^{+\infty}u_n(x)$ 在 $(a,b)$ 内闭一致收敛, 则 $\sum\limits_{n=1}^{+\infty}u_n(x)$ 的和函数在 $(a,b)$ 内连续。

我们前面已经知道 $\{x^n\}$ 在 $(-1,1)$ 内闭一致收敛, 从而它的极限函数 $0$ 在 $(-1,1)$ 中连续。

\noindent\textbf{例 10.4.1}\quad 设 $\{a_n\}$ 是一个单调序列且 $\lim\limits_{n\to\infty}a_n=0$, 证明级数 $\sum\limits_{n=1}^{+\infty}a_n\cos nx$ 与 $\sum\limits_{n=1}^{+\infty}a_n\sin nx$ 的和函数在 $(0,2\pi)$ 内连续。

\noindent\textbf{证明}\quad 对于 $\forall n\in\mathbb N$, 令
\[
  u_n(x)=a_n,\qquad v_n(x)=\cos nx.
\]
下面证明 $\sum\limits_{n=1}^{+\infty}u_n(x)v_n(x)$ 在 $(0,2\pi)$ 内闭一致收敛. 事实上, 任取闭区间 $[\delta_0,2\pi-\delta_0](\delta_0>0)$, 对于 $\forall n\in\mathbb N$, 我们有
\[
  \left|\sum_{k=1}^n\cos kx\right|
  =\frac{\left|\sin\left(n+\frac12\right)x-\sin\frac{x}{2}\right|}{2\sin\frac{x}{2}}
  \leq\frac1{\sin\frac{\delta_0}{2}}.
\]
由于 $u_n(x)$ 与 $x$ 无关, $\{u_n(x)\}$ 关于 $n$ 单调且在 $[\delta_0,2\pi-\delta_0]$ 上一致趋于零, 根据狄利克雷判别法, $\sum\limits_{n=1}^{+\infty}a_n\cos nx$ 在 $[\delta_0,2\pi-\delta_0]$ 上一致收敛。

由于 $a_n\cos nx\in C(-\infty,+\infty)(n=1,2,\cdots)$, 因此
\[
  \sum_{n=1}^{+\infty}a_n\cos nx\in C[\delta_0,2\pi-\delta_0].
\]
由 $\delta_0$ 的任意性, 我们有
\[
  \sum_{n=1}^{+\infty}a_n\cos nx\in C(0,2\pi).
\]
利用
\[
  \left|\sum_{k=1}^n\sin kx\right|
  =\frac{\left|\cos\left(n+\frac12\right)x-\cos\frac{x}{2}\right|}{2\sin\frac{x}{2}}
\]
同理可证 $\sum\limits_{n=1}^{+\infty}a_n\sin nx\in C(0,2\pi)$。

值得指出的是, 定理 10.4.1 及定理 10.4.1$'$ 都是给出了极限函数连续的充分条件. 容易举出例子来说明这些条件不是必要的, 例如 $f_n(x)=nxe^{-nx^2}(n=1,2,\cdots)$ 在 $[0,1]$ 点点收敛到连续函数 $0$, 但它在 $[0,1]$ 上不一致收敛; 又如函数项级数 $\sum\limits_{n=0}^{+\infty}(1-x^2)x^n=1+x, x\in(-1,1)$, 其和函数在 $(-1,1)$ 内连续, 但该函数项级数在 $(-1,1)$ 内不一致收敛。

下面的狄尼定理告诉我们, 如果附加一定的条件, 我们可以从点收敛推出一致收敛性。

\noindent\textbf{定理 10.4.2 (狄尼(Dini)定理)}\quad 设函数 $f_n(x)(n=1,2,\cdots)$ 在区间 $[a,b]$ 上连续, 且对于 $\forall x\in[a,b]$ 及 $\forall n\in\mathbb N$ 成立 $f_n(x)\leq f_{n+1}(x)$, 再设对于 $\forall x\in[a,b]$, 有 $\lim\limits_{n\to\infty}f_n(x)=f(x)$(即 $\{f_n(x)\}$ 是点点收敛的单调函数序列), 则 $f(x)$ 在区间 $[a,b]$ 上连续的充分必要条件是
\[
  f_n(x)\rightrightarrows f(x)\quad(x\in[a,b]).
\]

\noindent\textbf{证明}\quad 充分性是显然的, 现证必要性. 我们用反证法, 倘若 $\{f_n(x)\}$ 在 $[a,b]$ 上不一致收敛于 $f(x)$, 则 $\exists\varepsilon_0>0$, 对于 $\forall N\in\mathbb N$, 总存在 $n'>N$ 及 $x'\in[a,b]$, 使得
\[
  |f_{n'}(x')-f(x')|\geq\varepsilon_0.
\]
因此存在趋于无穷的正整数序列 $\{n_k\}$ 及序列 $\{x_{n_k}\}\subset[a,b]$, 使得
\[
  |f_{n_k}(x_{n_k})-f(x_{n_k})|\geq\varepsilon_0.
\]
我们不妨设 $x_{n_k}\to x_0\in[a,b](k\to\infty)$. 由 $f_n(x_0)\to f(x_0)(n\to\infty)$, 对于 $\varepsilon=\dfrac{\varepsilon_0}{4}$, 存在 $N\in\mathbb N$, 当 $n>N$ 时, 有
\[
  |f_n(x_0)-f(x_0)|<\varepsilon.
\]
取 $n_0=N+1$, 再由 $f_{n_0}(x)$ 及 $f(x)$ 的连续性, $\exists\delta>0$, 当 $x\in U(x_0,\delta)\cap[a,b]$ 时, 有
\[
  |f(x)-f(x_0)|<\varepsilon\quad\text{及}\quad |f_{n_0}(x)-f_{n_0}(x_0)|<\varepsilon.
\]
现在取一个 $n_k>n_0$, 使得 $|x_{n_k}-x_0|<\delta$, 则有
\[
  f_{n_0}(x_{n_k})\leq f_{n_k}(x_{n_k})\leq f(x_{n_k}).
\]
因此有
\[
\begin{aligned}
  \varepsilon_0
  &\leq |f_{n_k}(x_{n_k})-f(x_{n_k})|
  \leq |f_{n_0}(x_{n_k})-f(x_{n_k})|\\
  &\leq |f_{n_0}(x_{n_k})-f_{n_0}(x_0)|+|f_{n_0}(x_0)-f(x_0)|+|f(x_0)-f(x_{n_k})|\\
  &<3\varepsilon=\frac34\varepsilon_0.
\end{aligned}
\]
此矛盾便证明了必要性. 证毕。

由函数序列的狄尼定理, 我们容易推出下述关于函数项级数的狄尼定理。

\noindent\textbf{定理 10.4.2$'$ (狄尼定理)}\quad 设函数项级数 $\sum\limits_{n=1}^{+\infty}u_n(x)$ 在区间 $[a,b]$ 上收敛, 并且对于 $\forall n\in\mathbb N$, $u_n(x)$ 在 $[a,b]$ 上连续且非负, 则 $\sum\limits_{n=1}^{+\infty}u_n(x)$ 的和函数在 $[a,b]$ 上连续的充分必要条件是 $\sum\limits_{n=1}^{+\infty}u_n(x)$ 在 $[a,b]$ 上一致收敛。

\noindent\textbf{注}\quad 在狄尼定理中, 闭区间不能改为开区间或无穷区间, 请读者自己举例说明之。

\noindent\textbf{例 10.4.2}\quad 证明 $\sum\limits_{n=1}^{+\infty}x^n\dfrac{\ln x}{1+\left|\ln\ln\dfrac1x\right|}$ 在 $(0,1)$ 内一致收敛。

\noindent\textbf{证明}\quad 对于 $n=1,2,\cdots$, 令
\[
  u_n(x)=x^n\frac{\ln x}{1+\left|\ln\ln\dfrac1x\right|}\quad(x\in(0,1)).
\]
我们有 $\lim\limits_{x\to0+0}u_n(x)=0$ 和 $\lim\limits_{x\to1-0}u_n(x)=0$. 记
\[
  \widetilde u_n(x)=
  \begin{cases}
    0, & x=0,1,\\
    -u_n(x), & x\in(0,1),
  \end{cases}
\]
则 $\widetilde u_n(x)(n=1,2,\cdots)$ 在 $[0,1]$ 上连续、非负, 且
\[
  \sum_{k=1}^n\widetilde u_k(x)=
  \begin{cases}
    0, & x=0,1,\\
    -\dfrac{x-x^{n+1}}{1-x}\dfrac{\ln x}{1+\left|\ln\ln\dfrac1x\right|}, & x\in(0,1).
  \end{cases}
\]
因此我们有
\[
  \sum_{n=1}^{+\infty}\widetilde u_n(x)=
  \begin{cases}
    0, & x=0,1,\\
    -\dfrac{x}{1-x}\dfrac{\ln x}{1+\left|\ln\ln\dfrac1x\right|}, & x\in(0,1).
  \end{cases}
\]
容易看出它在 $[0,1]$ 上连续. 由定理 10.4.2$'$ 知 $\sum\limits_{n=1}^{+\infty}\widetilde u_n(x)$ 在 $[0,1]$ 一致收敛, 从而 $\sum\limits_{n=1}^{+\infty}u_n(x)$ 在 $(0,1)$ 内一致收敛, 即 $\sum\limits_{n=1}^{+\infty}x^n\dfrac{\ln x}{1+\left|\ln\ln\dfrac1x\right|}$ 在 $(0,1)$ 内一致收敛。

读者可以验证, 例 10.4.2 中的函数项级数不能用魏尔斯特拉斯 M 判别法来证明它是一致收敛的。

\subsection{极限函数的积分}

对于 $[a,b]$ 上可积函数序列 $\{f_n(x)\}$, 设它在 $[a,b]$ 上一致收敛到其极限函数 $f(x)$, 我们来研究 $f(x)$ 在 $[a,b]$ 上的可积性. 我们有以下的定理。

\noindent\textbf{定理 10.4.3}\quad 设函数 $f_n(x)(n=1,2,\cdots)$ 在 $[a,b]$ 上可积, 且 $f_n(x)\rightrightarrows f(x)(x\in[a,b])$, 则 $f(x)$ 在 $[a,b]$ 上可积, 并且成立
\[
  \lim_{n\to\infty}\int_a^b f_n(x)\mathrm dx
  =\int_a^b f(x)\mathrm dx
  =\int_a^b\lim_{n\to\infty}f_n(x)\mathrm dx.
\]

\noindent\textbf{证明}\quad 对于 $\forall\varepsilon>0$, 由 $f_n(x)\rightrightarrows f(x)(x\in[a,b])$, 存在 $N\in\mathbb N$, 当 $n>N$ 时, 对于 $\forall x\in[a,b]$, 有
\begin{equation}
  |f_n(x)-f(x)|<\frac{\varepsilon}{3(b-a)}.
  \tag{10.4.1}
\end{equation}
特别地, 取 $n_0=N+1$, 则 $f_{n_0}(x)$ 满足 (10.4.1) 式. 由 $f_{n_0}(x)$ 在 $[a,b]$ 上可积, 从而存在 $[a,b]$ 的分割
\[
  a=x_0<x_1<\cdots<x_n=b,
\]
若记 $\Delta x_i=x_i-x_{i-1}(i=1,2,\cdots,n)$, 则有
\[
  \sum_{i=1}^n \omega_i(f_{n_0})\Delta x_i<\frac\varepsilon3,
\]
其中我们用 $\omega_i(g)$ 表示函数 $g(x)$ 在 $[x_{i-1},x_i](i=1,2,\cdots,n)$ 上的振幅。

对 $x',x''\in[x_{i-1},x_i]$, 由 (10.4.1), 我们有
\[
\begin{aligned}
  |f(x')-f(x'')|
  &\leq |f(x')-f_{n_0}(x')|+|f_{n_0}(x')-f_{n_0}(x'')|+|f_{n_0}(x'')-f(x'')|.
\end{aligned}
\]
由上式我们推出
\[
  \omega_i(f)\leq\omega_i(f_{n_0})+\frac{2\varepsilon}{3(b-a)}.
\]
因此我们有
\[
\begin{aligned}
  \sum_{i=1}^n \omega_i(f)\Delta x_i
  &\leq \sum_{i=1}^n \omega_i(f_{n_0})\Delta x_i+\frac{2\varepsilon(b-a)}{3(b-a)}\\
  &\leq \frac\varepsilon3+\frac{2\varepsilon}{3}=\varepsilon.
\end{aligned}
\]
这就证明了 $f(x)$ 在 $[a,b]$ 上可积. 再由 (10.4.1), 当 $n>N$ 时, 有
\[
  \left|\int_a^b f_n(x)\mathrm dx-\int_a^b f(x)\mathrm dx\right|
  \leq \int_a^b |f_n(x)-f(x)|\mathrm dx
  \leq \frac{\varepsilon}{3(b-a)}(b-a)<\varepsilon,
\]
因此有
\[
  \lim_{n\to\infty}\int_a^b f_n(x)\mathrm dx=\int_a^b f(x)\mathrm dx.
\]
证毕。

同理可证下面的定理。

\noindent\textbf{定理 10.4.3$'$}\quad 设函数 $u_n(x)(n=1,2,\cdots)$ 在区间 $[a,b]$ 上可积, 且 $\sum\limits_{n=1}^{+\infty}u_n(x)$ 在 $[a,b]$ 上一致收敛, 则 $\sum\limits_{n=1}^{+\infty}u_n(x)$ 的和函数在 $[a,b]$ 上可积, 并且成立
\begin{equation}
  \int_a^b\left(\sum_{n=1}^{+\infty}u_n(x)\right)\mathrm dx
  =\sum_{n=1}^{+\infty}\int_a^b u_n(x)\mathrm dx.
  \tag{10.4.2}
\end{equation}

为了方便起见, 当 (10.4.2) 成立时, 我们也称函数项级数 $\sum\limits_{n=1}^{+\infty}u_n(x)$ 可在 $[a,b]$ 上逐项积分。

\noindent\textbf{例 10.4.3}\quad 证明: 当 $x\in(-1,1)$ 时, 成立以下等式:
\[
  \frac12\ln\frac{1+x}{1-x}
  =\sum_{n=0}^{+\infty}\frac{x^{2n+1}}{2n+1}.
\]

\noindent\textbf{证明}\quad 设 $f(x)=\dfrac1{1-x^2}$, 则当 $|x|<1$ 时, $f(x)$ 是以下几何级数的和函数:
\[
  f(x)=\frac1{1-x^2}=\sum_{n=0}^{+\infty}x^{2n}.
\]
由于其部分和序列
\[
  S_n(x)=\sum_{k=0}^n x^{2k}=\frac{1-x^{2n+2}}{1-x^2}\quad(n=1,2,\cdots)
\]
在 $(-1,1)$ 内闭一致收敛到 $\dfrac1{1-x^2}$, 因此对固定的 $x\in(-1,1)$, 该函数项级数可在 $[0,x](x>0)$ 或 $[x,0](x<0)$ 上逐项积分. 因此有
\[
  \frac12\ln\frac{1+x}{1-x}
  =\int_0^x\frac{\mathrm dt}{1-t^2}
  =\sum_{n=0}^{+\infty}\int_0^x t^{2n}\mathrm dt
  =\sum_{n=0}^{+\infty}\frac{x^{2n+1}}{2n+1}.
\]

\noindent\textbf{例 10.4.4}\quad 求 $I=\int_0^1\sum\limits_{n=1}^{+\infty}\dfrac{x}{n(n+x)}\mathrm dx$ 的值。

\noindent\textbf{解}\quad 对于 $\forall x\in[0,1]$ 及 $n=1,2,\cdots$, 有
\[
  0\leq\frac{x}{n(n+x)}\leq\frac1{n^2},
\]
因此 $\sum\limits_{n=1}^{+\infty}\dfrac{x}{n(n+x)}$ 在 $[0,1]$ 上一致收敛. 由定理 10.4.3$'$, 我们可以对该函数项级数逐项积分
\[
\begin{aligned}
  I
  &=\sum_{n=1}^{+\infty}\int_0^1\frac{x}{n(n+x)}\mathrm dx
   =\sum_{n=1}^{+\infty}\int_0^1\left(\frac1n-\frac1{n+x}\right)\mathrm dx\\
  &=\sum_{n=1}^{+\infty}\left[\frac1n-\ln(1+n)+\ln n\right]
   =\lim_{n\to\infty}\left[\sum_{k=1}^n\frac1k-\ln(1+n)\right]\\
  &=c\text{(欧拉常数)}.
\end{aligned}
\]

\subsection{极限函数的导数}

在本小节中, 我们研究函数序列（函数项级数）逐项求导组成的函数序列（函数项级数）的敛散性问题. 在这里我们注意到以下的一个事实: 即使一个函数序列 $\{f_n(x)\}$ 中每个函数在区间 $I$ 内都有连续的导函数 $f_n'(x)$, 且 $\{f_n'(x)\}$ 在区间 $I$ 上一致收敛, 也不能保证 $\{f_n(x)\}$ 在区间 $I$ 上一致收敛, 甚至不能保证 $\{f_n(x)\}$ 在区间 $I$ 上的点收敛性. 这是因为一个函数的原函数可相差一个常数的缘故. 要使得相应的 $\{f_n(x)\}$ 收敛, 则必须要求 $\{f_n(x)\}$ 至少在 $I$ 上的一个点上是收敛的. 我们有以下定理。

\noindent\textbf{定理 10.4.4}\quad 设函数 $f_n(x)(n=1,2,\cdots)$ 在区间 $[a,b]$ 上可微, 且满足

(a) 存在 $x_0\in[a,b]$, 使得 $\lim\limits_{n\to\infty}f_n(x_0)$ 存在;

(b) $f_n'(x)\rightrightarrows g(x)(x\in[a,b])$,

则我们有以下结论:

(1) 存在 $[a,b]$ 上的函数 $f(x)$, 使得 $f_n(x)\rightrightarrows f(x)(x\in[a,b])$;

(2) $f(x)$ 在 $[a,b]$ 上可微（端点单侧可微）, 并且 $f'(x)=g(x)$, 即
\[
  \lim_{n\to\infty}f_n'(x)=\left[\lim_{n\to\infty}f_n(x)\right]'.
\]

\noindent\textbf{证明}\quad 我们主要利用柯西准则来证明上述结论. 对于 $\forall\varepsilon>0$, 由 $\lim\limits_{n\to\infty}f_n(x_0)$ 存在和 $f_n'(x)\rightrightarrows g(x)(x\in[a,b])$ 知, $\exists N\in\mathbb N$, 当 $n>m\geq N$ 时, 有
\begin{equation}
  |f_n(x_0)-f_m(x_0)|<\frac\varepsilon2;
  \tag{10.4.3}
\end{equation}
并且对 $\forall x\in[a,b]$, 有
\begin{equation}
  |f_n'(x)-f_m'(x)|<\frac{\varepsilon}{2(b-a)}.
  \tag{10.4.4}
\end{equation}
对任意 $n>m\geq N$, 对 $f_n(x)-f_m(x)$ 在 $x_0$ 与 $x$ 之间应用拉格朗日微分中值定理知: 在 $x_0$ 与 $x$ 之间存在 $\xi$, 使得
\begin{equation}
\begin{aligned}
  &\left|[f_n(x)-f_m(x)]-[f_n(x_0)-f_m(x_0)]\right|\\
  &\qquad =|f_n'(\xi)-f_m'(\xi)||x-x_0|
   <\frac{|x-x_0|\varepsilon}{2(b-a)}\leq\frac\varepsilon2.
\end{aligned}
  \tag{10.4.5}
\end{equation}
因此对于 $\forall x\in[a,b]$, 由 (10.4.3) 和 (10.4.5) 得
\[
\begin{aligned}
  |f_n(x)-f_m(x)|
  &\leq \left|[f_n(x)-f_m(x)]-[f_n(x_0)-f_m(x_0)]\right|
      +|f_n(x_0)-f_m(x_0)|\\
  &\leq\frac\varepsilon2+\frac\varepsilon2=\varepsilon.
\end{aligned}
\]
由柯西准则知 $\{f_n(x)\}$ 在 $[a,b]$ 上一致收敛, 从而存在 $[a,b]$ 上的函数 $f(x)$, 使得 $f_n(x)\rightrightarrows f(x)(x\in[a,b])$, 即 (1) 得证。

为了证明 (2), 对于 $\forall x^*\in[a,b]$, 对 $n=1,2,\cdots$, 令
\[
  h_n(x)=\frac{f_n(x)-f_n(x^*)}{x-x^*}\ (x\ne x^*),\qquad h_n(x^*)=f_n'(x^*),
\]
\[
  h(x)=\frac{f(x)-f(x^*)}{x-x^*}\quad (x\ne x^*),
\]
则 $h_n(x)$ 在 $[a,b]$ 上连续. 现在我们证明 $\{h_n(x)\}$ 在 $[a,b]\setminus\{x^*\}$ 上一致收敛到 $h(x)$。

事实上, 从 (10.4.5) 的证明过程中可以推出, 对于 $\forall\varepsilon>0$, $\exists N\in\mathbb N$, 当 $n>m\geq N$ 时, 对于 $\forall x\in[a,b]\setminus\{x^*\}$, 有
\[
\begin{aligned}
  |h_n(x)-h_m(x)|
  &\leq \frac1{|x-x^*|}\left|[f_n(x)-f_m(x)]-[f_n(x^*)-f_m(x^*)]\right|\\
  &\leq\frac{\varepsilon}{2(b-a)}.
\end{aligned}
\]
因此 $h_n(x)$ 在 $[a,b]\setminus\{x^*\}$ 上一致收敛到 $h(x)$. 注意到 $\lim\limits_{x\to x^*}h_n(x)=f_n'(x^*)$, 由定理 10.4.1 后的注, 我们有 $f'(x^*)=\lim\limits_{x\to x^*}h(x)$ 存在, 并且 $f'(x^*)=\lim\limits_{n\to\infty}f_n'(x^*)$. 由 $x^*\in[a,b]$ 的任意性, 因此 (2) 成立. 证毕。

对于函数项级数, 我们有

\noindent\textbf{定理 10.4.4$'$}\quad 设函数 $u_n(x)(n=1,2,\cdots)$ 在区间 $[a,b]$ 上可微, 且满足:

(a) 存在 $x_0\in[a,b]$, 使得 $\sum\limits_{n=1}^{+\infty}u_n(x_0)$ 收敛;

(b) $\sum\limits_{n=1}^{+\infty}u_n'(x)$ 在 $[a,b]$ 上一致收敛,

则

(1) $\sum\limits_{n=1}^{+\infty}u_n(x)$ 在 $[a,b]$ 上一致收敛;

(2) $\sum\limits_{n=1}^{+\infty}u_n(x)$ 的和函数在 $[a,b]$ 上可导, 并且
\begin{equation}
  \left(\sum_{n=1}^{+\infty}u_n(x)\right)'
  =\sum_{n=1}^{+\infty}u_n'(x).
  \tag{10.4.6}
\end{equation}

同样为了便于叙述, 当 (10.4.6) 成立时, 我们称函数项级数 $\sum\limits_{n=1}^{+\infty}u_n(x)$ 可在 $[a,b]$ 上逐项求导数（或逐项可微）。

\noindent\textbf{注}\quad 由于函数可微也是一种局部性质, 因此要讨论函数序列的极限函数的可微性只需要求该函数序列点点收敛, 其导函数序列内闭一致收敛即可. 显然, 对于函数项级数也有类似的结论。

\noindent\textbf{例 10.4.5}\quad 证明 $\sum\limits_{n=1}^{+\infty}\dfrac1{n^x}$ 在 $(1,+\infty)$ 上无穷次可微。

\noindent\textbf{证明}\quad 任取 $\delta_0>0$, 我们已经知道 $\sum\limits_{n=1}^{+\infty}\dfrac1{n^x}$ 在 $(1+\delta_0,+\infty)$ 上一致收敛. 由于
\[
  \left(\frac1{n^x}\right)'=\frac{-\ln n}{n^x}\in C(0,+\infty),
\]
并且对于 $\forall x\in(1+\delta_0,+\infty)$, 有
\[
  \left|\frac{-\ln n}{n^x}\right|<\frac{\ln n}{n^{1+\delta_0}},
\]
注意到 $\sum\limits_{n=1}^{+\infty}\dfrac{\ln n}{n^{1+\delta_0}}$ 收敛, 因此 $\sum\limits_{n=1}^{+\infty}\dfrac{-\ln n}{n^x}$ 在 $(1+\delta_0,+\infty)$ 上一致收敛。

由定理 10.4.4$'$ 知 $\sum\limits_{n=1}^{+\infty}\dfrac1{n^x}$ 的和函数在 $(1+\delta_0,+\infty)$ 上有连续的导函数, 且
\[
  \left(\sum_{n=1}^{+\infty}\frac1{n^x}\right)'
  =\sum_{n=1}^{+\infty}\frac{-\ln n}{n^x}.
\]
由于 $\sum\limits_{n=1}^{+\infty}\dfrac{-\ln n}{n^x}$ 逐项求导后得到的函数项级数 $\sum\limits_{n=1}^{+\infty}\dfrac{\ln^2 n}{n^x}$ 仍然在 $(1+\delta_0,+\infty)$ 上一致收敛, 所以 $\sum\limits_{n=1}^{+\infty}\dfrac1{n^x}$ 的和函数在 $(1+\delta_0,+\infty)$ 上有连续的二阶导数. 依次下去知, $\sum\limits_{n=1}^{+\infty}\dfrac1{n^x}$ 的和函数在 $(1+\delta_0,+\infty)$ 上具有任意阶导数. 再由 $\delta_0$ 的任意性, 知 $\sum\limits_{n=1}^{+\infty}\dfrac1{n^x}$ 的和函数在 $(1,+\infty)$ 上无穷次可微。

利用函数项级数, 我们可以构造许多具有特殊性质的函数。

\noindent\textbf{例 10.4.6}\quad 设 $\{x_n\}\subset[0,1]$ 是一个各项互不相同的序列, 证明 $\sum\limits_{n=1}^{+\infty}\dfrac{\operatorname{sgn}(x-x_n)}{2^n}$ 的和函数当且仅当 $x=x_n(n=1,2,\cdots)$ 时不连续。

\noindent\textbf{证明}\quad 首先注意到, 对任意的 $n$ 及 $x\in[0,1]$, 有
\[
  \left|\frac{\operatorname{sgn}(x-x_n)}{2^n}\right|\leq\frac1{2^n}.
\]
由于 $\sum\limits_{n=1}^{+\infty}\dfrac1{2^n}$ 收敛, 因此 $\sum\limits_{n=1}^{+\infty}\dfrac{\operatorname{sgn}(x-x_n)}{2^n}$ 在 $[0,1]$ 上一致收敛. 对于 $\forall k\in\mathbb N$, 该函数项级数中仅有一项 $\dfrac{\operatorname{sgn}(x-x_k)}{2^k}$ 在 $x=x_k$ 处不连续, 而其他项均在 $x=x_k$ 处连续. 因此 $\sum\limits_{n=1}^{+\infty}\dfrac{\operatorname{sgn}(x-x_n)}{2^n}$ 的和函数在 $x=x_k$ 处不连续。

当 $x_0\notin\{x_n\}$ 时, 由于该函数项级数每一项都在 $x=x_0$ 处连续, 从而 $\sum\limits_{n=1}^{+\infty}\dfrac{\operatorname{sgn}(x-x_n)}{2^n}$ 的和函数在 $x=x_0$ 处连续. 证毕。

\noindent\textbf{注}\quad 若在上例中, 将符号函数 $\operatorname{sgn}(x)$ 换为
\[
  f(x)=
  \begin{cases}
    x^2\sin\dfrac1x, & x\ne0,\\
    0, & x=0,
  \end{cases}
\]
作
\[
  g(x)=\sum_{n=1}^{+\infty}\frac{f(x-x_n)}{2^n},
\]
则可证明 $g(x)$ 在 $[0,1]$ 上可导, 但 $g'(x)$ 在 $x=x_n\in[0,1](n=1,2,\cdots)$ 处不连续, 而在 $[0,1]$ 中的其他点连续。

历史上有许多人曾经以为一个连续函数仅有可能在一列点上是不可导的. 魏尔斯特拉斯利用函数项级数首先构造出了一个处处不可导的连续函数. 后来, 人们借助函数项级数构造了一些更为简单的处处不可导的连续函数的例子, 如 V. D. Waerden 构造了以下例子:

记函数 $\widetilde f(x)=|x|, x\in\left[-\dfrac12,\dfrac12\right]$, 然后记将 $\widetilde f(x)$ 以 $1$ 为周期延拓到 $(-\infty,+\infty)$ 上的函数为 $f(x)$, 则 $\sum\limits_{n=1}^{+\infty}\dfrac{f(4^n x)}{4^n}$ 的和函数是 $(-\infty,+\infty)$ 上的一个连续函数, 但它在 $(-\infty,+\infty)$ 上处处不可导。

由于这个事实的证明较为复杂, 我们在这里省略之。

\section*{习题十}

1. 求出下列函数项级数的收敛域:
\[
  (1)\ \sum_{n=1}^{+\infty}\frac{n}{n+1}\left(\frac{x}{2x+1}\right)^n;\qquad
  (2)\ \sum_{n=1}^{+\infty}\frac{x^n}{1-x^n}.
\]

2. 求下列函数序列 $\{f_n(x)\}$ 的极限函数, 并讨论在给定的区间上是否一致收敛:
\[
\begin{gathered}
  (1)\ f_n(x)=x(1-x)^n,\ x\in[0,1].\\
  (2)\ f_n(x)=nxe^{-nx^2},\ x\in(-\infty,+\infty).\\
  (3)\ f_n(x)=n^2e^{-nx^2},\ x\in[a,+\infty),\ \text{这里 }a>0.\\
  (4)\ f_n(x)=\frac{x^2}{x^2+(1-nx)^2},\ x\in[0,1].\\
  (5)\ f_n(x)=n^2x(1-x^2)^n,\ x\in[0,1].\\
  (6)\ f_n(x)=\sin\frac{x}{n^n};\ \text{(a) }x\in[a,b];\ \text{(b) }x\in(-\infty,+\infty).\\
  (7)\ f_n(x)=\frac{\sin n^n x}{n^\alpha},\ \alpha>0,\ x\in(-\infty,+\infty).\\
  (8)\ f_n(x)=(\sin x)^{n^\alpha},\ \alpha>0,\ x\in[0,\pi].
\end{gathered}
\]

3. 讨论下列函数序列或函数项级数在指定区间的一致收敛性。
\[
\begin{gathered}
  (1)\ \left\{f_n(x)=x^{\sum_{k=1}^n\frac1{2^k}}\right\},\ x\in(0,1);\\
  (2)\ \sum_{n=1}^{+\infty}\frac1{n^2+x^2},\ x\in(-\infty,+\infty);\\
  (3)\ \sum_{n=1}^{+\infty}\frac{\sin x\sin nx}{\sqrt{n+x^2}},\ x\in[0,+\infty);\\
  (4)\ \sum_{n=1}^{+\infty}\frac{(-1)^n}{n+x^2},\ x\in(-\infty,+\infty);
\end{gathered}
\]

\end{document}
